\documentclass[conference]{IEEEtran}
\usepackage{amsmath,amssymb}
\usepackage{booktabs}
\usepackage{url}
\usepackage{hyperref}

\title{Extending a BPMN Discrete-Event Simulator for Resource-Aware Cache-System Studies}
\author{Author Name\\Department / Institution\\Email address}

\begin{document}
\maketitle

\begin{abstract}
This paper describes a plugin-oriented extension to a BPMN discrete-event simulator for studying distributed
cache-service behavior. The extension adds a resource-aware service-time model based on CPU, storage I/O and network
demands, a generalized-mean (softmax) aggregation, CPU idle-capacity scaling, and hard and soft affinity constraints
that can refer to resource properties or to the execution history of the same process instance. A small resource-tuple
ranking hook integrates affinity with the existing resource scheduler. Three synthetic experiments illustrate (i)
cache-tier scale-out in the presence of a fixed database bottleneck, (ii) fixed network-transfer penalties combined
with ownership and replica constraints, and (iii) heterogeneous cache nodes under uniform routing and soft-affinity
ranking. The scenarios are provided as BPMN and Scylla configuration files. The current model does not maintain cache
contents or a network topology; those limits are made explicit so that the experiments are interpreted as controlled
what-if studies rather than as simulations of a complete cache implementation.
\end{abstract}

\noindent\textbf{Keywords---} BPMN, discrete-event simulation, distributed cache, resource-aware service time,
affinity scheduling, softmax

\section{Introduction}

Cache systems are useful examples for evaluating interactions between request routing, heterogeneous resources,
replication and backend bottlenecks. A BPMN model expresses the request path and branching behavior, while a
discrete-event engine represents request arrivals, resource queues and task execution. A basic process simulator,
however, usually treats task duration as an input distribution and does not directly connect task work to the selected
machine or expose execution history to the resource scheduler.

This work describes an extension that preserves the simulator's BPMN execution and resource-queue machinery while
adding these capabilities through plugins. A single hook in the core scheduler allows extensions to rank candidate
resource tuples after the scheduler has checked resource type, quantity and availability. Separate parser and event
plugins interpret hardware/workload annotations, calculate task service times and record per-process-instance
resource history.

The contribution is an implementable, reproducible modeling layer rather than a complete distributed-cache engine.
Three generated experiments show how the extension can be used and identify which conclusions are supported by the
current abstractions.

\section{Extension Design}

\subsection{Resource-Aware Task Demand}

A workload is attached to a task in the simulation configuration. The supported demand attributes are instruction
count $N_{\mathrm{instr}}$, read and write operation counts $N_{\mathrm{read}}$ and $N_{\mathrm{write}}$, and transmitted
and received data volumes $D_{\mathrm{tx}}$ and $D_{\mathrm{rx}}$. Hardware attributes are attached to a dynamic
resource type or a named resource instance: CPU speed $S_{\mathrm{cpu}}$, idle factor $\alpha$, read/write rates
$B_{\mathrm{read}}$ and $B_{\mathrm{write}}$, and transmit/receive rates $B_{\mathrm{tx}}$ and $B_{\mathrm{rx}}$.
Instance attributes override type defaults. The parser plugins retain these values in the simulator's extension data;
no resource-model field or core XML parser is changed.

CPU service time includes idle-capacity scaling:
\begin{equation}
S_{\mathrm{cpu}}^{\mathrm{eff}} = S_{\mathrm{cpu}}(1-\alpha), \qquad
T_{\mathrm{cpu}} = \frac{N_{\mathrm{instr}}}{S_{\mathrm{cpu}}(1-\alpha)}, \quad 0\leq\alpha<1.
\end{equation}
Storage and network components are additive within their domains:
\begin{align}
T_{\mathrm{io}} &= \frac{N_{\mathrm{read}}}{B_{\mathrm{read}}}+
                    \frac{N_{\mathrm{write}}}{B_{\mathrm{write}}},\\
T_{\mathrm{net}} &= \frac{D_{\mathrm{tx}}}{B_{\mathrm{tx}}}+
                     \frac{D_{\mathrm{rx}}}{B_{\mathrm{rx}}}.
\end{align}
The total service time combines the three components with a generalized mean:
\begin{equation}
T_{\mathrm{service}} = \left(T_{\mathrm{cpu}}^p+T_{\mathrm{io}}^p+T_{\mathrm{net}}^p\right)^{1/p},
\qquad p\geq 1.
\label{eq:softmax}
\end{equation}
For $p=1$, the result is additive; increasing $p$ approaches the largest component. The implementation normalizes by
the largest component before exponentiation to reduce overflow for large $p$. Durations are expressed in seconds when
speeds are expressed per second. The simulator represents event time at whole-second resolution and truncates
fractional task durations, so example values are scaled to avoid relying on sub-second differences.

When more than one resource is assigned to a task, the current duration plugin uses the lowest defined speed for each
component as the bottleneck. A component with nonzero demand but no corresponding speed is reported as an error. This
is a modeling convention; it is not a fluid-sharing or parallel-service model.

\subsection{Affinity-Aware Resource Selection}

The core scheduler first forms candidate resource tuples that satisfy requested resource types, amounts and timetables,
then applies the existing default last-access ordering. A new \texttt{ResourceTupleRankingPluggable} hook allows a
plugin to filter and reorder these candidates before assignment. The affinity plugin implements two kinds of rules:

\begin{itemize}
  \item \emph{Hard rules} reject candidates that violate affinity or anti-affinity.
  \item \emph{Soft rules} add their configured weights to candidates that satisfy them; candidates are sorted by the
        sum of satisfied weights, preserving default order for ties.
\end{itemize}

Rules may target a resource type or named instance, or refer to the resources that executed an earlier task of the
same process instance. If a history reference has not executed yet, its rule does not apply. History is recorded at
task start and is scoped to the process instance. Ranking is side-effect-free because the scheduler also requests a
candidate selection when it is checking whether an event can proceed.

If all candidates are rejected, the task remains queued. In that case the core returns the resources temporarily
polled during candidate construction to their queues. This preserves queue state and allows a later retry. The ranking
hook itself does not reserve resources; the existing scheduler performs the assignment after ranking.

A soft weight is a ranking priority, not a probability. With the current plugin, a weight of four does not imply four
times as many assignments as a weight of one. Similarly, a rule-free task is allocated by the simulator's default
resource policy, not necessarily uniform random or round robin.

\subsection{Gateway Seeding}

The existing simulator seeds distributions in a configuration with the same configured seed. Consequently, an
arrival-rate distribution and an exclusive-gateway distribution can produce correlated random sequences. An opt-in
system property, \texttt{-Dscylla.independentSeeds=true}, derives an exclusive-gateway seed from the configuration
seed and gateway node identifier. The default behavior is retained for compatibility. The cache experiment runner
enables the option when comparing results with independent Poisson thinning; it does not change seeds for all other
distribution types.

\section{Scenario Models}

All experiments are synthetic and use 20,000 process instances, mean inter-arrival time 50 simulated seconds and a
fixed 0.9 hit probability unless specified otherwise. The illustrative time mapping is one simulated second to
100 ms. Speeds are abstract parameters, not measurements of hardware. The BPMN gateway fixes the hit/miss fraction;
there is no cache-content state and no key-popularity-driven hit calculation.

\subsection{Scale-Out with a Fixed Database Tier}

The process is the distributed-cache flow: cache lookup, exclusive hit/miss split, return on hit, and database read,
cache write and replication on miss. Cache nodes are identical, with $S_{\mathrm{cpu}}=1000$ instructions/s and
$B_{\mathrm{tx}}=1$ MB/s. A single database resource has $B_{\mathrm{read}}=1$ operation/s. The four configurations
vary only the cache population $N\in\{2,4,8,16\}$; no affinity rules are present.

A lookup consumes 10,000 instructions and transmits 10 MB, hence 10 s CPU and 10 s network service at $p=1$. Cache
write and replication each consume 20,000 instructions (20 s). A database read consumes 200 read operations (200 s).
With arrival rate $\lambda=1/50$ and miss probability $q=0.1$, the nominal database utilization is
\begin{equation}
\rho_{\mathrm{DB}}=\lambda q T_{\mathrm{DB}}=\frac{1}{50}\cdot 0.1\cdot 200=0.4.
\end{equation}
Under ideal even allocation, approximate total cache demand is $20+q(20+20)=24$ seconds per request, yielding
per-node utilization $\rho_{\mathrm{cache}}\approx24/(50N)$. This gives 0.24, 0.12, 0.06 and 0.03 for the four
cluster sizes. These values are analytical estimates, not simulation measurements. The expected signature of a fixed
backend bottleneck is falling cache-node utilization with little improvement in miss latency once DB waiting/service
dominates.

Measurements should include per-node utilization and assignment share, database utilization and queue wait, mean hit
and miss response times, and completion time for the fixed workload. Averages alone may conceal imbalance, so both
per-node and aggregate utilization should be reported.

\subsection{Fixed Transfer-Volume Locality Proxy}

The second experiment uses four equal cache instances. The intended labels are Zone A for \#0 and \#1 and Zone B for
\#2 and \#3. Six simulation configurations combine three fixed transfer-volume categories with two assignment
variants:
\begin{center}
\begin{tabular}{lrr}
\toprule
Configured path category & $D_{\mathrm{tx}}$ (MB) & $D_{\mathrm{tx}}/B_{\mathrm{tx}}$ at 1 MB/s\\
\midrule
local & 0 & 0 s\\
same zone & 2 & 2 s\\
cross zone & 40 & 40 s\\
\bottomrule
\end{tabular}
\end{center}
In the \texttt{none} variants, lookup has no affinity rule and the simulator's default resource policy applies. In the
\texttt{owner} variants, lookup is hard-pinned to instance \#0; cache write has hard affinity to the resource recorded
for lookup, and replication has hard anti-affinity to that resource.

These configurations are a controlled proxy for locality costs, not a topology-aware network simulation. The value
$D_{\mathrm{tx}}$ is the task's fixed data-volume demand and is independent of which node is selected. Therefore the
same-zone and cross-zone service costs do not arise from node-pair distance, and \#0 is treated as the owner in every
\texttt{owner} configuration. The model has no zone labels, key-owner map, per-link bandwidth, propagation latency or
link contention. A topology-aware version would need a plugin/model extension that determines the selected node and
owner, then computes a per-pair network cost.

\subsection{Heterogeneous Nodes and Routing}

The third experiment assigns three cache instances with equal network speed $B_{\mathrm{tx}}=1$ MB/s but different
CPU speeds: $S_{\mathrm{cpu},0}=2000$, $S_{\mathrm{cpu},1}=1000$, and $S_{\mathrm{cpu},2}=500$ instructions/s. A
lookup workload of 10,000 instructions and 10 MB has nominal $p=1$ service times 15, 20 and 30 s on nodes \#0, \#1
and \#2.

The uniform-routing control uses a separate BPMN with an exclusive gateway whose branches each have probability 1/3.
Each branch leads to a task hard-pinned to one cache instance. It provides probabilistic uniform assignment, not
deterministic round robin; if the selected node is busy, the request waits for that node rather than failing over.

The comparison configuration uses the distributed-cache BPMN and gives the single cache-lookup task soft-affinity
weights 4, 2 and 1 for the three instances. The current ranker selects the highest-scoring available candidate; it
does not allocate in a 4:2:1 probability ratio. The comparison therefore evaluates fast-node preference with
availability-based fallback, not proportional weighted routing. A true proportional or capacity-aware policy
requires weighted random assignment or a scheduler that explicitly considers expected service time.

The principal metrics are lookup assignment fraction by node, per-node utilization, queueing delay at each node, and
hit/miss response time. The uniform baseline is expected to route one third of requests to each node, while soft
ranking should favor the faster node when it is among the feasible candidates. The simulation configs should be run
with independent gateway seeds when an independent routing/arrival comparison is intended.

\section{Reproducibility}

The generated files are under \path{samples/other_cache_scenarios/}. Regenerate them from the repository root with:
\begin{verbatim}
python3 samples/cache_scenarios/run_experiments.py write-other-scenarios
\end{verbatim}
The configurations have not been executed for this paper draft; the outcomes above are analytical expectations and
hypotheses to test, not empirical results. For reproducible runs, use the same global seed across paired variants,
record the software revision, enable BPS/XES logging, and use
\texttt{-Dscylla.independentSeeds=true} when random gateway decisions should be independent of arrivals. Store each
run's inputs and output log separately.

\section{Limitations and Future Work}

The current scenarios do not compute cache hits from contents, model key skew, derive zones or owners from topology,
model network propagation latency, or dynamically route in proportion to soft-affinity weights. The service-time model
uses fixed per-resource speeds and does not share CPU, I/O or link capacity continuously among concurrent tasks.
Whole-second truncation also limits studies of small delays. Future extensions should add per-node cache state and
replacement policies, an explicit network graph with pair-specific latency/bandwidth, and a policy interface for
weighted or capacity-aware placement. Each addition should be evaluated against an analytically tractable baseline
before interpreting larger-scale results.

\section{Conclusion}

The extension connects BPMN task workloads to heterogeneous resource capabilities and makes affinity/history visible
to resource selection while retaining the simulator's existing event and queue model. The three cache scenarios
provide controlled experiments for scale-out, transfer-cost sensitivity and heterogeneous-node routing. Their most
important methodological result is also a boundary: fixed workload and ranking rules can demonstrate scheduler
behavior, but they do not by themselves implement cache state, topology-dependent latency or probabilistic weighted
routing.

\end{document}
